
\title{The Japanese DP and the structure of its modifiers}
\author{Viktor Köhlich}

\BackBody{This book offers an analysis of the Japanese DP in line with the cartographic research program. It addresses theoretical linguists interested in nominal modification as well as scholars in the field of Japanese studies. The main results concern the identification of three domains, a DP-external domain, a domain for indirect modifiers and a domain for direct modifiers, with specific discussion on the latter two, and an account for the internal structure of nominal modifiers. On the one hand, this book highlights the applicability of the cartographic program to languages outside the Indo-European cosmos. At the same time, it offers a re-evaluation of many previously made assumptions based on the Japanese data, eventually allowing for refinements of these assumptions and a better understanding of the DP in cross-linguistic perspective.}
\renewcommand{\lsSeries}{ogs}
\renewcommand{\lsSeriesNumber}{14}
\renewcommand{\lsID}{588}
\typesetter{Viktor Köhlich, Sebastian Nordhoff}
\proofreader{Akiko Nagano,
Brett Reynolds,
David Carrasco Coquillat,
Elliott Pearl,
Ikmi Nur Oktavianti,
Mary Ann Walter,
Matthew Korte,
Nicoletta Romeo,
Jean Nitzke,
Wilson Lui
}
\renewcommand{\lsImpressumExtra}{Arbeit zur Erlangung des Doktorgrades (Sigelziffer D.30)}
\renewcommand{\lsISBNdigital}{978-3-96110-575-5}
\renewcommand{\lsISBNhardcover}{978-3-98554-196-6}
\BookDOI{10.5281/zenodo.20410026}
